\documentclass[12pt,a4paper]{article}
\usepackage[utf8]{inputenc}
\usepackage[T1]{fontenc}
\usepackage[polish]{babel}
\usepackage{geometry}
\usepackage{graphicx}
\usepackage{xcolor}
\usepackage{listings}
\usepackage{float}
\usepackage{caption}
\usepackage{hyperref}
\usepackage{amsmath}
\geometry{margin=2.5cm}

\lstset{
  basicstyle=\ttfamily\small,
  breaklines=true,
  frame=single,
  tabsize=2,
  showstringspaces=false,
  keywordstyle=\color{blue!60!black},
  commentstyle=\color{green!40!black},
  stringstyle=\color{orange!50!black}
}

\title{Sprawozdanie z implementacji oświetlenia i odbić w silniku graficznym}
\author{Jędrzej Gołaszewski \\ Album: 331690}
\date{\today}

\begin{document}
\maketitle
\tableofcontents
\newpage

\section{Cel i zakres pracy}
Celem sprawozdania jest analiza implementacji oświetlenia w projekcie silnika graficznego, ze szczególnym uwzględnieniem:
\begin{itemize}
  \item elementów oświetlenia sceny (światła kierunkowe/pozycyjne),
  \item wykorzystania światła odbitego,
  \item wpływu materiału na odbicia powierzchniowe oraz rozproszone,
  \item roli cieni i ich filtracji w finalnym obrazie.
\end{itemize}
Zakres obejmuje część merytoryczną oraz istotne fragmenty kodu źródłowego projektu.

\section{Część merytoryczna}
\subsection{Model oświetlenia zastosowany w projekcie}
W shaderze pikseli zastosowano klasyczny model oparty o składowe:
\begin{itemize}
  \item \textbf{ambient} -- stałe oświetlenie otoczenia,
  \item \textbf{diffuse (Lambert)} -- oświetlenie rozproszone zależne od kąta padania światła,
  \item \textbf{specular} -- odbicie powierzchniowe, modelowane przez potęgowanie składowej oświetlenia.
\end{itemize}
Składowe są sumowane dla każdej lampy, a następnie uzupełniane o emissive materiału. W praktyce profesjonalnej taki układ (ambient + diffuse + specular) jest często punktem wyjścia w silnikach czasu rzeczywistego, ponieważ zapewnia dobrą czytelność obrazu przy umiarkowanym koszcie obliczeniowym. Dzięki temu model sprawdza się w etapach prototypowania oraz jako baza do późniejszego przejścia na bardziej zaawansowane BRDF.

W zapisie matematycznym można to przedstawić jako:
\[
I = I_{a} + \sum_{i=1}^{N}\left(I_{d,i} + I_{s,i}\right) + I_{e}
\]
gdzie $I_a$ to składowa ambient, $I_{d,i}$ i $I_{s,i}$ to odpowiednio diffuse i specular dla $i$-tego światła, a $I_e$ oznacza emissive.

Dla modelu Lambert:
\[
I_{d,i} = k_d\,L_i\,\max(\mathbf{N}\cdot\mathbf{L}_i, 0)
\]
natomiast składowa powierzchniowa (w uproszczeniu zgodnym z implementacją):
\[
I_{s,i} = k_s\,L_i\,\big(\max(\mathbf{N}\cdot\mathbf{L}_i, 0)\big)^p
\]
gdzie $p$ jest potęgą połysku zależną od chropowatości.

\subsection{Wpływ materiału na odbicia i rozproszenie}
Parametry materiału przekazywane do shadera:
\begin{itemize}
  \item \texttt{BaseColor} -- kolor bazowy,
  \item \texttt{Emissive} -- samorozświetlenie,
  \item \texttt{Metallic} -- udział charakteru metalicznego,
  \item \texttt{Roughness} -- chropowatość powierzchni.
\end{itemize}
Dla wyższego \texttt{Metallic} rośnie udział odbicia powierzchniowego (specular), natomiast \texttt{Roughness} wpływa na ostrość/refleksyjność odbić poprzez zmianę parametru mocy specular. Taki zestaw parametrów jest zgodny z profesjonalnym workflow materiałowym opartym o PBR (Physically Based Rendering), ponieważ pozwala utrzymać spójny wygląd obiektów przy różnych konfiguracjach świateł. Dodatkowo ograniczanie wartości do sensownych zakresów stabilizuje rendering i zmniejsza ryzyko artefaktów podczas pracy zespołowej (programista + grafik techniczny).

W projekcie kolor rozproszony i powierzchniowy jest mieszany analogicznie do poniższych zależności:
\[
\mathbf{c}_{diff} = \mathbf{BaseColor}\,(1 - m)
\]
\[
\mathbf{c}_{spec} = (1-m)\,\mathbf{F}_0 + m\,\mathbf{BaseColor}, \quad \mathbf{F}_0 \approx (0.04,0.04,0.04)
\]
gdzie $m \in [0,1]$ oznacza \texttt{Metallic}. Dla chropowatości $r$ przyjęto:
\[
g = 1-r, \quad p = \operatorname{lerp}(8,128,g)
\]
co odpowiada przejściu od szerokiego, matowego refleksu do ostrzejszego połysku.

\subsection{Odbicia powierzchniowe i rozproszone}
\begin{itemize}
  \item \textbf{Odbicie rozproszone} dominuje dla materiałów mniej metalicznych i bardziej chropowatych.
  \item \textbf{Odbicie powierzchniowe} staje się silniejsze dla materiałów metalicznych i gładszych.
\end{itemize}
Zależność ta jest kluczowa dla realizmu odwzorowania różnych typów powierzchni. W profesjonalnej grafice komputerowej właśnie kontrola proporcji między składową rozproszoną i powierzchniową decyduje o wiarygodności materiału (np. plastik, ceramika, metal). W docelowych pipeline'ach produkcyjnych ten sam kierunek jest rozwijany przez bardziej fizyczne modele BRDF (np. GGX/Cook-Torrance), ale logika interpretacji materiału pozostaje analogiczna.

Udział energii między składowymi można zapisać schematycznie jako:
\[
I_{final} \propto \alpha\,I_{diff} + \beta\,I_{spec}, \quad \alpha \approx (1-m),\; \beta \approx m
\]
co pokazuje, że wzrost metaliczności zwiększa wagę odbicia powierzchniowego względem rozproszonego.

\subsection{Cienie i ich wpływ na percepcję materiału}
W projekcie użyto map cieni z filtrowaniem PCF (3x3), co ogranicza aliasing krawędzi cienia. Dodatkowo wykorzystano bias oraz normal bias, aby zmniejszyć artefakty typu self-shadowing. W kontekście profesjonalnym jest to często stosowany kompromis jakości do kosztu: miększe krawędzie cieni poprawiają percepcję głębi i kontaktu obiektów z podłożem, a jednocześnie rozwiązanie pozostaje wystarczająco lekkie dla aplikacji czasu rzeczywistego. Z tego powodu podobne podejście jest szeroko używane w grach i interaktywnych wizualizacjach.

Dla pojedynczej próbki test głębokości można zapisać jako:
\[
V = H\big(z_{map}(u,v) - (z_{frag}-b)\big)
\]
gdzie $H$ to funkcja skokowa, $b$ to bias, $z_{frag}$ to głębokość fragmentu w przestrzeni światła. W przypadku PCF 3\(\times\)3 końcowa widoczność ma postać:
\[
V_{PCF} = \frac{1}{9}\sum_{x=-1}^{1}\sum_{y=-1}^{1} V_{x,y}
\]
co odpowiada uśrednieniu dziewięciu porównań i daje łagodniejsze przejście na granicy cienia.

\section{Istotne fragmenty kodu z projektu}
\subsection{Shader pikseli -- obliczenia diffuse/specular}
\begin{lstlisting}[language=C++,caption={GraphicEngineApp/shaders/color_ps.hlsl -- fragment głównej pętli oświetlenia}]
float3 N = normalize(input.normalWS);
float metallic = saturate(g_EmissiveMetallic.w);
float roughness = saturate(g_SurfaceParams.x);

float3 diffuseColor = g_BaseColor.rgb * (1.0f - metallic);
float3 specularColor = lerp(float3(0.04f, 0.04f, 0.04f), g_BaseColor.rgb, metallic);
float gloss = 1.0f - roughness;
float specularPower = lerp(8.0f, 128.0f, gloss);

float3 diffuseAccum = 0.0f;
float3 specAccum = 0.0f;

[loop]
for(uint i = 0; i < g_LightCount; ++i)
{
    float3 lightDir = normalize(g_Lights[i].DirectionIntensity.xyz);
    float intensity = g_Lights[i].DirectionIntensity.w;
    float3 lightColor = g_Lights[i].ColorUsePosition.rgb;
    float usePosition = g_Lights[i].ColorUsePosition.w;

    float3 L = normalize(-lightDir);
    float attenuation = 1.0f;

    float lambert = max(dot(N, L), 0.0f);
    float specular = pow(lambert, specularPower);
    float shadow = SampleShadowPCF(i, input.worldPos, N, L);

    diffuseAccum += diffuseColor * lightColor * (lambert * intensity * attenuation * shadow);
    specAccum += specularColor * lightColor * (specular * intensity * attenuation * shadow);
}

float3 ambient = 0.04f * diffuseColor;
float3 emissive = g_EmissiveMetallic.rgb;

return float4(ambient + diffuseAccum + specAccum + emissive, g_BaseColor.a);
\end{lstlisting}

\subsection{Shader pikseli -- cienie (PCF)}
\begin{lstlisting}[language=C++,caption={GraphicEngineApp/shaders/color_ps.hlsl -- funkcja SampleShadowPCF}]
float SampleShadowPCF(uint lightIndex, float3 worldPos, float3 normalWS, float3 lightVec)
{
    if(g_ObjectParams.x < 0.5f)
    {
        return 1.0f;
    }

    const LightData light = g_Lights[lightIndex];
    const float shadowEnabled = light.ShadowParams.w;
    const float shadowMapIndex = light.ShadowParams.x;
    if(shadowEnabled < 0.5f || shadowMapIndex < 0.0f)
    {
        return 1.0f;
    }

    float4 lightClip = mul(float4(worldPos, 1.0f), light.LightViewProj);
    if(abs(lightClip.w) < 1e-5f)
    {
        return 1.0f;
    }

    float3 lightNdc = lightClip.xyz / lightClip.w;
    if(lightNdc.x < -1.0f || lightNdc.x > 1.0f || lightNdc.y < -1.0f || lightNdc.y > 1.0f || lightNdc.z < 0.0f || lightNdc.z > 1.0f)
    {
        return 1.0f;
    }

    uint shadowWidth;
    uint shadowHeight;
    uint shadowElements;
    uint shadowLevels;
    g_ShadowMaps.GetDimensions(0, shadowWidth, shadowHeight, shadowElements, shadowLevels);

    float2 uv = lightNdc.xy * float2(0.5f, -0.5f) + 0.5f;
    float2 texelSize = 1.0f / float2(max(shadowWidth, 1), max(shadowHeight, 1));

    const float NdotL = saturate(dot(normalWS, lightVec));
    const float bias = light.ShadowParams.y + (1.0f - NdotL) * light.ShadowParams.z;
    const float compareDepth = saturate(lightNdc.z - bias);

    const float slice = shadowMapIndex;
    float visibility = 0.0f;

    [unroll]
    for(int y = -1; y <= 1; ++y)
    {
        [unroll]
        for(int x = -1; x <= 1; ++x)
        {
            float2 sampleUv = uv + float2((float)x, (float)y) * texelSize;
            visibility += g_ShadowMaps.SampleCmpLevelZero(g_ShadowSampler, float3(sampleUv, slice), compareDepth);
        }
    }

    return visibility / 9.0f;
}
\end{lstlisting}

\subsection{Materiał -- ustawianie właściwości}
\begin{lstlisting}[language=C++,caption={GraphicEngineCore/src/renderer/Material.cpp -- SetProperties}]
void Material::SetProperties(const MaterialProperties& properties, nvrhi::ICommandList* commandList) {
    m_Constants.BaseColor = properties.BaseColor;
    m_Constants.EmissiveMetallic = {
        properties.Emissive.x,
        properties.Emissive.y,
        properties.Emissive.z,
        std::clamp(properties.Metallic, 0.0f, 1.0f)
    };
    m_Constants.SurfaceParams = {
        std::clamp(properties.Roughness, 0.0f, 1.0f),
        0.0f,
        0.0f,
        0.0f
    };

    UploadConstants(commandList);
}

void Material::SetRoughness(float roughness, nvrhi::ICommandList* commandList) {
    m_Constants.SurfaceParams.x = std::clamp(roughness, 0.0f, 1.0f);
    UploadConstants(commandList);
}
\end{lstlisting}

\subsection{Tworzenie obiektów testowych o różnych parametrach materiału}
\begin{lstlisting}[language=C++,caption={GraphicEngineApp/SandboxApp.cpp -- konfiguracja Metallic/Roughness}]
m_Cubes.push_back(
    Renderer::EntityBuilder(device, cmd)
    .Cube()
    .Material(m_DefaultVS, m_DefaultPS)
    .Position(0.0f, 0.0f, 0.0f)
    .Scale(2.0f, 2.0f, 2.0f)
    .Color(0.95f, 0.95f, 0.95f, 1.0f)
    .Emissive(0.0f, 0.0f, 0.0f)
    .Metallic(0.0f)
    .Roughness(0.8f)
    .Build()
);

m_Cubes.push_back(
    Renderer::EntityBuilder(device, cmd)
    .Sphere()
    .Material(m_DefaultVS, m_DefaultPS)
    .Position(0.0f, -5.0f, 5.0f)
    .Scale(2.0f, 2.0f, 2.0f)
    .Color(0.95f, 0.95f, 0.95f, 1.0f)
    .Emissive(0.0f, 0.0f, 0.0f)
    .Metallic(0.8f)
    .Roughness(0.8f)
    .Build()
);
\end{lstlisting}

\section{Metodyka testów aplikacji}
Wszystkie testy należy wykonywać przy porównywalnych warunkach (ta sama scena, podobna pozycja kamery, kontrolowane parametry światła). Każdy test powinien zawierać:
\begin{itemize}
  \item cel testu,
  \item zestaw parametrów wejściowych,
  \item obserwację jakościową,
  \item zrzut ekranu dokumentujący wynik.
\end{itemize}

\section{Sterowanie i interfejs użytkownika}
\subsection{Sterowanie kamerą i ruchem w scenie}
W aplikacji zastosowano sterowanie zbliżone do standardu używanego w narzędziach 3D i silnikach czasu rzeczywistego. Ruch kamery realizowany jest klawiszami:
\begin{itemize}
  \item \texttt{W}/\texttt{S} -- ruch przód/tył,
  \item \texttt{A}/\texttt{D} -- ruch lewo/prawo,
  \item \texttt{Q}/\texttt{E} -- ruch dół/góra.
\end{itemize}
Obrót kamery odbywa się po przytrzymaniu prawego przycisku myszy, co pozwala na płynną inspekcję sceny i ocenę wpływu materiałów oraz oświetlenia z różnych kierunków.

\textbf{Miejsce na zdjęcie (widok sceny podczas sterowania kamerą):}
\begin{figure}[H]
\centering
\fbox{\rule{0pt}{6cm}\rule{0.9\linewidth}{0pt}}
\caption{Sterowanie kamerą -- przykładowy widok}
\end{figure}

\subsection{Interfejs użytkownika i kontrola parametrów}
Interfejs użytkownika (oparty o ImGui) wspiera iteracyjną pracę nad sceną: umożliwia podgląd i bieżące strojenie parametrów renderingu. W kontekście profesjonalnym taki panel jest istotny, ponieważ skraca czas testowania ustawień świateł i materiałów bez potrzeby każdorazowej rekompilacji kodu.

W szczególności interfejs może służyć do:
\begin{itemize}
  \item obserwacji aktualnego stanu sceny i parametrów,
  \item szybkiej zmiany ustawień wpływających na wygląd obrazu,
  \item porównywania wyników między kolejnymi konfiguracjami testowymi.
\end{itemize}

\textbf{Miejsce na zdjęcie (panel interfejsu użytkownika):}
\begin{figure}[H]
\centering
\fbox{\rule{0pt}{6cm}\rule{0.9\linewidth}{0pt}}
\caption{Interfejs użytkownika -- panel kontroli}
\end{figure}

\subsection{Znaczenie ergonomii obsługi}
Czytelne sterowanie oraz prosty panel UI bezpośrednio wpływają na jakość analizy wyników. Dzięki wygodnej nawigacji i szybkiemu dostępowi do parametrów można szybciej wykrywać różnice w odbiciach powierzchniowych i rozproszonych, a także rzetelniej dokumentować testy.

\textbf{Miejsce na zdjęcie (porównanie dwóch konfiguracji):}
\begin{figure}[H]
\centering
\fbox{\rule{0pt}{6cm}\rule{0.9\linewidth}{0pt}}
\caption{Porównanie konfiguracji A/B}
\end{figure}

\section{Testy aplikacji i miejsce na dokumentację zdjęciową}
\subsection{Test 1: Wpływ parametru Roughness (Metallic = const)}
\textbf{Parametry:}
\begin{itemize}
  \item Metallic: 0.0 (materiał niemetaliczny)
  \item Roughness: 0.1, 0.5, 0.9 (trzy porównywane przebiegi)
  \item Ustawienia światła: światło główne białe, intensywność 7.2, stała pozycja kamery
\end{itemize}
\textbf{Miejsce na zdjęcie:}
\begin{figure}[H]
\centering
\fbox{\rule{0pt}{6cm}\rule{0.9\linewidth}{0pt}}
\caption{Test 1 -- screenshot}
\end{figure}
\textbf{Obserwacje:} Dla niskiego roughness pojawia się wyraźniejszy połysk i ostrzejsze refleksy. Przy roughness równym 0.9 powierzchnia staje się matowa, a składowa rozproszona wizualnie dominuje nad powierzchniową.

\subsection{Test 2: Wpływ parametru Metallic (Roughness = const)}
\textbf{Parametry:}
\begin{itemize}
  \item Metallic: 0.0, 0.5, 1.0 (trzy porównywane przebiegi)
  \item Roughness: 0.4 (wartość stała)
  \item Ustawienia światła: dwa światła spot, ten sam kadr i ta sama odległość kamery
\end{itemize}
\textbf{Miejsce na zdjęcie:}
\begin{figure}[H]
\centering
\fbox{\rule{0pt}{6cm}\rule{0.9\linewidth}{0pt}}
\caption{Test 2 -- screenshot}
\end{figure}
\textbf{Obserwacje:} Wraz ze wzrostem metallic maleje udział koloru rozproszonego, a rośnie wpływ odbicia powierzchniowego. Dla metallic bliskiego 1.0 materiał wizualnie przyjmuje charakter metalu, z wyraźnym odbiciem koloru światła.

\subsection{Test 3: Zmiana koloru i intensywności światła}
\textbf{Parametry:}
\begin{itemize}
  \item Kolor światła: ciepły (1.0, 0.92, 0.82) vs chłodny (0.45, 0.55, 1.0)
  \item Intensywność: 2.0, 4.0, 7.2
  \item Pozycja/kierunek: stała pozycja, kierunek na środek sceny
\end{itemize}
\textbf{Miejsce na zdjęcie:}
\begin{figure}[H]
\centering
\fbox{\rule{0pt}{6cm}\rule{0.9\linewidth}{0pt}}
\caption{Test 3 -- screenshot}
\end{figure}
\textbf{Obserwacje:} Zmiana barwy źródła bezpośrednio wpływa na odbiór materiałów i atmosferę sceny. Wyższa intensywność poprawia czytelność odbić, ale przy zbyt dużej wartości może prowadzić do prześwietlenia najjaśniejszych fragmentów.

\subsection{Test 4: Cienie włączone vs wyłączone}
\textbf{Parametry:}
\begin{itemize}
  \item CastShadows: OFF (porównanie bazowe) oraz ON (tryb docelowy)
  \item Shadow Bias: 0.0015
  \item Normal Bias: 0.0025
\end{itemize}
\textbf{Miejsce na zdjęcie:}
\begin{figure}[H]
\centering
\fbox{\rule{0pt}{6cm}\rule{0.9\linewidth}{0pt}}
\caption{Test 4 -- screenshot}
\end{figure}
\textbf{Obserwacje:} Po włączeniu cieni wzrasta wrażenie głębi i kontaktu obiektów z podłożem. Ustawione bias i normal bias ograniczają artefakty typu acne, zachowując jednocześnie stabilność konturu cienia.

\subsection{Test 5: Wpływ zasięgu i kąta stożka światła spot}
\textbf{Parametry:}
\begin{itemize}
  \item Range: 8.0, 14.0, 20.0
  \item Inner/Outer angle: (20\degree/32\degree), (30\degree/46\degree)
  \item Pozycja źródła: test A -- blisko obiektu, test B -- wyżej i dalej od sceny
\end{itemize}
\textbf{Miejsce na zdjęcie:}
\begin{figure}[H]
\centering
\fbox{\rule{0pt}{6cm}\rule{0.9\linewidth}{0pt}}
\caption{Test 5 -- screenshot}
\end{figure}
\textbf{Obserwacje:} Większy zasięg zwiększa obszar oświetlenia, ale obniża lokalny kontrast. Szerszy stożek daje bardziej równomierne pokrycie sceny, natomiast węższy lepiej akcentuje wybrane detale materiału.

\section{Wnioski}
Na podstawie przeprowadzonych eksperymentów można stwierdzić, że parametry materiałowe \texttt{Metallic} i \texttt{Roughness} mają kluczowy wpływ na charakter odbić oraz sposób rozpraszania światła. Zastosowanie map cieni z filtrowaniem PCF znacząco poprawia czytelność geometrii i realizm sceny. W praktyce najlepsze efekty wizualne uzyskuje się przez równoważenie ustawień materiału i źródeł światła, z uwzględnieniem rodzaju powierzchni oraz zamierzonego efektu artystycznego.

\end{document}
